\section{A degree-two and degree-three stop for graph compatibility alone}
\label{sec:136-low-tate-compatibility-obstruction}

\paragraph{Archive, inputs, and exact scope.}
This concerns cote 136, physical pp.~117, 119, 121, footers
116, 118, 120, scans 117, 119, 121.  The question tested here
is whether the common-graph primitive property for all old
degree-three torsion pairs is preserved by a free degree-two
or degree-three Tate stage.  The answer is negative, already
in the category truncated in simplicial dimensions zero, one,
and two.  The old algebra has integer-torsion second cohomology,
not vanishing second cohomology.  This distinction is essential.

The construction uses the checked supplied first-skeleton
replacement, its free left Kan envelope, and the canonical
relative comparison and integral simplex-period normalization
from \nolinkurl{ARCH-136-SKELETON-OBSTRUCTION}.  These are
inherited internal inputs.  The low-degree algebra, the
all-modulus primitive calculation, and the two coordinate
changes are proved here (outside inputs: none; no A, M, or
outside preprint).  Reviewer123 proposed and checked the mixed
residue formula motivating the example and the degree-three
coordinate change.  The parent proposed the degree-two branch.
Their subsequent whole-proof readings, if completed, are recorded
separately from those contributions.

Let \(S=\Gamma_{\mathbb Z}(T)\), and let \(B\) be the
standard integral theory.  Use the checked first-skeleton
replacement to fix an admissible theory and a strict map
\[
 \rho:C\longrightarrow B,\qquad C_0=S,\quad C_1=F_1.
                                                        \tag{136.LTQ1}
\]
Here \(F_1\) is the actual supplied free algebra of even
zero-form unit disks, with its specified coefficient endpoints.
Write \(L=L(F)\) for its left Kan envelope.  It has a
canonical map to \(C\), and each \(L_n\) is a free
tensor of even zero-form unit disks over \(S\).  In
particular it contracts \(S\)-linearly to the coefficients.
Only its restriction to simplicial dimensions at most two will
be needed below.

Fix a prime \(p\) and positive weights \(r,s\), chosen
as large as any prescribed finite target guards.  Put
\(w=r+s\).  The relative comparison supplies an actual form
\[
 \eta\in Z^2 C^{w}(\Delta[2],\partial\Delta[2]),
          \qquad \int_{\Delta[2]}\rho(\eta)=T^{[w]}.
                                                        \tag{136.LTQ2}
\]
Indeed relative second cohomology is canonically \(S^w\)
in this range, and \(\rho\) preserves that coefficient
class.  The standard integral divided-power integration is
the canonical map with this normalization.  Set
\(G=\partial\Delta[2]=\operatorname{sk}_1\Delta[2]\).

\paragraph{ARCH-136-Q2-Q3-COMPATIBILITY-OBSTRUCTION.}
There is a graded-free, rationally vertically exact,
two-truncated source \(A\rightrightarrows C\) with literal
first skeleton \(F\), with \(H^0(A_2)=S\),
\(H^1(A_2)=0\), and integer-torsion \(H^2(A_2)\),
having the following properties.
Every old pair \((e,\xi)\in A_2^2\oplus A_2^3\), with
\[
 |\xi|=3,\quad |e|=2,\qquad d\xi=0,\quad de=N\xi,
                 \qquad N\geq1,
                                                        \tag{136.LTQ3}
\]
has target primitives \(c_j,t_j\) satisfying
\(dc_j=f_j\xi\), \(dt_j=Nc_j-f_je\), with both
restrictions to \(G\) the same for \(j=0,1\).
Nevertheless each of the following two free, full-boundary
relative two-cell repairs has target extensions, and every
choice of the extensions has nonzero degree-three
graph-difference residue in arbitrarily high weights:
a degree-three Tate pair and a degree-two Tate pair.
Both repaired sources still have \(H^0=S\) and
\(H^1=0\) in every weight and simplex through dimension
two, so the supplied immediate RH1 repair is the identity
on this truncation.  Thus RH1 does not remove either failure.

This is a failure of preservation for prescribed old maps.
The two target theories are the same admissible \(C\),
and hence already have a common source.  No counterexample to
the archive connectedness or homotopy question is asserted.
No all-pair compatibility assertion in higher simplicial
dimensions is implicit in this two-truncated theorem.

\begin{proof}
Use the free special differential PD algebra \(D\) over
\(S\) with odd generators \(y,z,a\), even generators
\(x,b,k\), and data
\[
 \begin{array}{c|rrrrrr}
        &y&x&b&z&a&k\\ \hline
 \text{form}&1&2&2&3&3&2\\
 \text{weight}&r&r&s&s&w&w
 \end{array}
\]
\[
 \begin{aligned}
 dy&=p x,& dx&=0,\\
 db&=p z,& dz&=0,\\
 da&=x b-y z,& dk&=p^2a-p yb.
 \end{aligned}
                                                        \tag{136.LTQ4}
\]
All differentials square to zero.  For example
\(d(xb-yz)=p xz-p xz=0\), and
\(d(p^2a-pyb)=p^2(xb-yz)-p(px b-pyz)=0\).
The special derivation extends to the even divided powers;
the derivative of each even generator is odd and has square
zero by strict graded commutativity.

Define \(A_0=S\), \(A_1=F_1\), and
\[
                         A_2=L_2\otimes_S D.
                                                        \tag{136.LTQ5}
\]
Every new generator has zero on every face of the two-cell.
The faces use the augmentation \(D\to S\), and all
degeneracies from lower simplices are those of \(L\).
This defines an actual two-truncated simplicial differential
PD algebra: the nonlinear expressions in (136.LTQ4) also
have zero faces.  Define two maps to the corresponding
truncation of \(C\), canonical on \(L\), by
\[
 \begin{aligned}
 f_0(y)&=f_0(x)=f_0(b)=f_0(z)=f_0(a)=f_0(k)=0,\\
 f_1(k)&=p\eta,\quad
 f_1(y)=f_1(x)=f_1(b)=f_1(z)=f_1(a)=0.
 \end{aligned}
                                                        \tag{136.LTQ6}
\]
The equations and faces are respected, so these are strict
maps and agree literally on the supplied first skeleton.

We first compute all old cohomology needed, including the
relations, over the torsion-free coefficient ring \(S\).
In form one there is \(S y\); in form two the basis is
\(x,b,k\); in form three the basis is
\(z,a,yx,yb,yk\).  Put
\[
                              h=p a-yb.
\]
The differentials of the last three basis elements are
\[
 \begin{aligned}
 d(yx)&=2p\gamma_2(x),\\
 d(yb)&=p(xb-yz),\\
 d(yk)&=p xk-p^2ya.
 \end{aligned}
                                                        \tag{136.LTQ7}
\]
The independent form-four basis monomials force the kernel
in form three to be exactly \(S z\oplus S h\).
The form-two boundaries into it are \(pS z\oplus pS h\).
Likewise the closed form-two module is \(Sx\), with
boundaries \(pSx\).  Thus
\[
 \begin{aligned}
 H^0(A_2)&=S,& H^1(A_2)&=0,\\
 H^2(A_2)&=(S/p)[x],&
 H^3(A_2)&=(S/p)[z]\oplus(S/p)[h].
 \end{aligned}
                                                        \tag{136.LTQ8}
\]
Here the calculation passes from \(D\) to \(A_2\)
by the actual \(S\)-linear contraction of \(L_2\),
so no tensor flatness inference is required.  Rationally
replace \(a\) by \(a'=a-yb/p\).  Then \(da'=0\)
and \(dk=p^2a'\), while the pairs \((y,x)\) and
\((b,z)\) have differential coefficient \(p\).
They are free unit disks after rational rescaling.  This
proves rational vertical exactness in every weight.  The
source is graded-free integrally and has literal structural
weight zero.

Here is a direct verification of the old all-modulus property;
in particular it is not inferred merely from checking a chosen
list of order-\(p\) relations.  For (136.LTQ3), (136.LTQ8)
gives homogeneous coefficients \(\lambda,\mu\in S\)
and an old two-form \(\alpha\) such that
\[
                         \xi=\lambda z+\mu h+d\alpha.
\]
Since \(N[\xi]=0\), there are coefficients
\(R,U\in S\) with \(N\lambda=pR\),
\(N\mu=pU\).  Subtracting a known primitive gives
a closed two-form; the second group in (136.LTQ8) then gives
an old one-form \(\beta\) and coefficient \(\ell\)
with
\[
                   e=Rb+Uk+N\alpha+\ell x+d\beta.
                                                        \tag{136.LTQ9}
\]
All equations respect weights; missing negative-weight
coefficient groups are zero.  Write \(\varepsilon_0=0\),
\(\varepsilon_1=1\).  The actual target choices
\[
 c_j=f_j\alpha+\varepsilon_j\mu\eta,
                      \qquad t_j=-f_j\beta
                                                        \tag{136.LTQ10}
\]
satisfy \(dc_j=f_j\xi\) and
\(dt_j=Nc_j-f_je\), by (136.LTQ6), (136.LTQ9),
and \(N\mu=pU\).  Their graph restrictions agree:
old restrictions agree and \(\eta|_G=0\).  This proves
the asserted property for every old pair and every modulus.

For the degree-three branch adjoin a free pair on the full
relative two-cell, of weight \(s\):
\[
 \begin{aligned}
 |c|&=2,\quad |t|=1,\qquad dc=z,\quad dt=pc-b,\\
 c|_G&=t|_G=0.
 \end{aligned}
                                                        \tag{136.LTQ11}
\]
Both old maps extend, for example by sending \(c,t\) to
zero.  Put
\[
 X=a+xt-yc,\qquad K=k+pyt,\qquad dX=0,\quad dK=p^2X.
                                                        \tag{136.LTQ12}
\]
These identities follow directly from (136.LTQ4) and
(136.LTQ11).  More strongly, set \(b_*=pc-b\).
The change to the free coordinates
\(y,x,t,b_*,c,z,K,X\) is triangular and invertible
as a strict graded PD change of variables.  In these coordinates
the only differentials are
\[
               dy=px,\qquad dt=b_*,\qquad dc=z,\qquad dK=p^2X.
                                                        \tag{136.LTQ13}
\]
Consequently \(H^0=S\), \(H^1=0\),
\(H^2=(S/p)[x]\), and \(H^3=(S/p^2)[X]\)
in simplex two.  To check the last assertion without any
K\"unneth shortcut, the even unit disk \((c,z)\)
contracts.  The arity-one block of the odd unit disk
\((t,b_*)\) contracts, and its remaining blocks begin
in form three with injective differential on coefficient
zeroforms.  After removing those unit blocks the remaining form-three
basis is \(X,yx,yK\), with differentials zero,
\(2p\gamma_2(x)\), and \(p xK-p^2yX\), respectively.
The last two are independent nonzero free-basis expressions.
Thus the third kernel is \(SX\), with boundaries
\(p^2SX\), as asserted.  The remaining higher blocks of
\((t,b_*)\) contribute no third cycle because their first
coefficient differential is injective over \(S\).
In particular \([X]\) has exact order \(p^2\).

For the degree-two branch, instead adjoin a free full-boundary
pair of weight \(r\):
\[
 \begin{aligned}
 |u|&=1,\quad |q|=0,\qquad du=x,\quad dq=pu-y,\\
 u|_G&=q|_G=0.
 \end{aligned}
                                                        \tag{136.LTQ14}
\]
Again zero images extend both old maps.  Define
\[
 a_0=a-ub-qz,\qquad k'=k-pqb,\qquad v=y-pu.
                                                        \tag{136.LTQ15}
\]
This is another triangular invertible free PD coordinate
change, now including an even zeroform.  All products and
translations are permitted by the free special-PD universal
property.  The differential becomes
\[
                    dq=-v,\qquad du=x,\qquad db=pz,\qquad dk'=p^2a_0.
                                                        \tag{136.LTQ16}
\]
In particular \(da_0=0\).  Contract the even zero-form
unit disk \((q,v)\) and the arity-one part of the odd
unit disk \((u,x)\).  The same low block calculation gives
\[
 \begin{aligned}
 H^0&=S,\quad H^1=H^2=0,\\
 H^3&=(S/p)[z]\oplus(S/p^2)[a_0].
 \end{aligned}
                                                        \tag{136.LTQ17}
\]
This also proves that the class of \(a_0\) has exact
order \(p^2\).  Both branches retain \(H^0=S\) and \(H^1=0\) in
lower simplices, since their first skeleton is unchanged.  RH1 is therefore the identity on
the full two-truncation: every normalized one-cycle already
has its unique zero-form primitive with the prescribed vertex
value.  This is the actual supplied RH1, not a rational
replacement.

It remains to check the obstruction for arbitrary retained
target extensions.  In the degree-three branch every such
extension satisfies
\[
                      f_jX=0,\qquad f_jK=\varepsilon_j p\eta,
                                                        \tag{136.LTQ18}
\]
because the old \(x,y,a\) have zero images.  In the
degree-two branch exactly the same equations hold with
\((X,K)\) replaced by \((a_0,k')\), because the
old \(b,z,a\) have zero images.  These facts do not
depend on the new generator images.  In either case denote
the resulting Moore pair by \((Z,H)\), so
\(dH=p^2 Z\), \(f_j Z=0\), and
\(f_jH=\varepsilon_jp\eta\).

If target primitives with common graph restrictions existed,
there would be forms \(A_j\) and \(Q_j\) of forms two
and one such that
\(dA_j=0\), \(dQ_j=p^2 A_j-f_jH\), with
\(A_1|_G=A_0|_G\), \(Q_1|_G=Q_0|_G\).
Apply \(\rho\), integrate the difference, and use Stokes:
\[
 0=\int_{\partial\Delta[2]}\rho(Q_1-Q_0)
   =p^2\int_{\Delta[2]}\rho(A_1-A_0)-pT^{[w]}.
                                                        \tag{136.LTQ19}
\]
The integral in the last expression is an element of
\(S^w\), by the standard integral divided-power period
formula.  Equation (136.LTQ19) is impossible because
\(T^{[w]}\) is a basis of the free rank-one group
\(S^w\).  Equivalently the graph-difference linking value
is \(T^{[w]}/p\) in \(\mathbb Q[T]/S\), up to the
single common sign convention.

Finally multiply this actual pair by \(T^{[p^a]}\),
where \(p^a>w\).  Its linking value is
\[
 \frac{1}{p}\binom{w+p^a}{w}\,T^{[w+p^a]}\pmod S.
                                                        \tag{136.LTQ20}
\]
The binomial coefficient is a unit modulo \(p\): in
\(\mathbb F_p[X]\),
\((1+X)^{w+p^a}=(1+X)^w(1+X^{p^a})\), and the
coefficient of \(X^w\) is one because \(w<p^a\).
Thus these are nonzero in arbitrarily high weights.  This
proves every stated failure, including its fixed-target and
post-RH1 scope.
\end{proof}

\paragraph{A variant with the minimal order in the last old Tate pair.}
The failure does not require the extra factor \(p\) in
the last old equation.  Keep every generator and differential
in (136.LTQ4), except replace that equation and its target image by
\[
                         dk=pa-yb=h,\qquad f_1(k)=\eta.
                                                        \tag{136.LTQ21}
\]
This is an actual minimal-order degree-four Tate pair over the
preceding algebra on \(y,x,b,z\).  Indeed
\(\theta=xb-yz\) is a four-cycle there, with
\[
 Z^4=S\gamma_2(x)\oplus S\theta,\qquad
 B^4=2pS\gamma_2(x)\oplus pS\theta.
                                                        \tag{136.LTQ22}
\]
For this calculation the form-four basis is
\(\gamma_2(x),xb,\gamma_2(b),yz\), whose derivatives
are respectively \(0,p xz,p bz,p xz\).  The
form-three derivatives are exactly \(0,2p\gamma_2(x),
p\theta\).  Thus \([\theta]\) has exact order
\(p\), and \(da=\theta\), \(dk=pa-yb\) uses
that order.

The old calculation now gives \(H^2=(S/p)[x]\) and
\(H^3=(S/p)[z]\), because the entire old \(S h\)
is a boundary.  For an arbitrary old pair, write
\(\xi=\lambda z+d\alpha\), \(N\lambda=pR\),
and \(e=Rb+N\alpha+\ell x+d\beta\).
The choices \(c_j=f_j\alpha\), \(t_j=-f_j\beta\)
again give common graph primitives for every modulus.  No
condition is lost by removing the old \([h]\) summand.

In the degree-three branch keep \(X=a+xt-yc\) but set
\(K=k+yt\).  Then \(dK=pX\).  The same invertible
coordinate change gives (136.LTQ13) with its last coefficient
\(p\), so \(H^3=(S/p)[X]\).  In the degree-two
branch keep \(a_0=a-ub-qz\) but put \(k'=k-qb\).
Then \(dk'=pa_0\), and the same coordinates give
\(H^3=(S/p)[z]\oplus(S/p)[a_0]\).  The low
\(H^0,H^1\) and RH1 assertions are unchanged.
Every target extension sends the indicated three-cycle to zero
and its relation to \(\varepsilon_j\eta\).  Stokes
would therefore require
\(p\int\rho(A_1-A_0)=T^{[w]}\), which is again
impossible.  The same coefficient multiples have nonzero
linking value (136.LTQ20).

This refinement only checks the minimal order of the displayed
last old degree-four pair.  It does not assert a complete
minimal construction history for the preceding scaled disks,
nor does it make the two-truncated source admissible.

\paragraph{ARCH-136-LOW-TATE-RETAINED-COMPLETION.}
The original diagram (136.LTQ4)--(136.LTQ6) cannot factor
through an admissible intermediate whose relative comparison
is in range at the two indicated weights.  More precisely,
let \(W_E(d)\) be a bound above which the canonical
relative comparison for an admissible \(E\) holds through
form degree \(d\), for every simplicial pair.  There is
no factorization of the given two-truncated maps by strict
coefficient-preserving arrows as
\[
 \begin{gathered}
 A\longrightarrow E_{\leq2}\rightrightarrows C_{\leq2},\\
 r\geq W_E(2),\qquad w\geq W_E(3).
 \end{gathered}
                                                        \tag{136.LTQ23}
\]
No agreement of the two maps on the whole first skeleton of
\(E\) is assumed.  The fixed target comparison bounds
for \(C\) are already part of the choice of \(r,s,w\).

\begin{proof}
All old positive generators have zero boundary, so their images
are relative forms for \((\Delta[2],G)\) in \(E\).
Its relative second cohomology in weight \(r\) is
canonically the free group \(S^r\).  Since \(dy=px\),
the relative class of \(x\) is zero.  Choose an actual
relative one-form \(u\) with \(du=x\).  The relative
one-form \(pu-y\) is closed, and relative first cohomology
vanishes, so choose a relative zeroform \(q\) with
\(dq=pu-y\).  These choices use only the first inequality
in (136.LTQ23).

Inside \(E\) define the same \(a_0=a-ub-qz\) and
\(k'=k-pqb\).  They are relative forms, with
\(da_0=0\) and \(dk'=p^2a_0\).  Relative third
cohomology in weight \(w\) is zero, so there is a relative
two-form \(P\) with \(dP=a_0\).  Therefore
\[
                             k'-p^2P
                                                        \tag{136.LTQ24}
\]
is a closed relative two-form.  The two strict coefficient maps
from \(E\) to \(C\) induce the same canonical map
\(S^w\to S^w\) on this relative cohomology group.
But each image of \(a_0\) is zero and the two images
of \(k'\) are \(0,p\eta\), just as before, independently
of the chosen \(u,q\).  Comparing the periods of
(136.LTQ24) therefore gives
\[
 pT^{[w]}=p^2\left(
       \int\rho f_1P-\int\rho f_0P\right).
                                                        \tag{136.LTQ25}
\]
Both integrals are integral elements of \(S^w\), which
is impossible.  This proves the guarded assertion using
relative comparison and naturality, without a common first
skeleton for the intermediate.
\end{proof}

In fact a coefficient shift eliminates both guards for the
\emph{single fixed old block}.  This refinement was proposed by
reviewer123.  Given any purported admissible intermediate
\(E\), choose \(p^a>w\) so large that
\[
 r+p^a\geq W_E(2),\qquad w+p^a\geq W_E(3),
                 \qquad \sigma=T^{[p^a]}.
                                                        \tag{136.LTQ26}
\]
The preceding relative argument, now applied to
\(\sigma x,\sigma y\), supplies relative forms
\(du=\sigma x\), \(dq=pu-\sigma y\).  Define
\[
 a_\sigma=\sigma a-ub-qz,\qquad
 k_\sigma=\sigma k-pqb.
                                                        \tag{136.LTQ27}
\]
Every term has weight \(w+p^a\).  Direct differentiation
using the old equations gives
\(da_\sigma=0\), \(dk_\sigma=p^2a_\sigma\).
No separate comparison guard on the old \(b,z\) is
needed: they enter only these products.  The two target images
of \(a_\sigma\) are zero, and those of
\(k_\sigma\) are \(0,p\sigma\eta\), for the
same reason as in (136.LTQ18).

Choose a relative \(P\) with \(dP=a_\sigma\)
using (136.LTQ26).  The canonical equality of the images
of the closed relative two-form \(k_\sigma-p^2P\)
would imply
\[
 p\binom{w+p^a}{w}T^{[w+p^a]}
     =p^2\left(\int\rho f_1P-\int\rho f_0P\right).
                                                        \tag{136.LTQ28}
\]
The binomial coefficient is a unit modulo \(p\) by
(136.LTQ20), and the two integrals are integral.  This is
impossible.  Thus no admissible retained intermediate exists
for this one fixed two-truncated diagram, with no bounds on
that intermediate left in the conclusion.  The old all-modulus
common-graph property proved in (136.LTQ9)--(136.LTQ10)
continues to hold for this very same source.

For completeness this also gives an actual full simplicial
source with the stated noncompletion property.  Left Kan
extend the fixed two-truncated algebra along the fully faithful
inclusion of the truncated simplex category.  The unit retains
its entire two-truncation, and its two maps to \(C\)
extend by the defining adjunction.  Any factorization through
an admissible theory would restrict to the excluded factorization
of that two-truncation.  No assertion about all old pairs in
higher simplices is required.  The full source is inadmissible,
since otherwise it would be an admissible intermediate itself.
Both targets are still the identical admissible \(C\).
This is noncompletion for specified retained data, not failure
of common-source existence or of homotopy for arrows with an
admissible domain.

\paragraph{What this closes and what remains.}
The mixed old-\(H^2\) piece of the low arity filtration
in Section~\ref{sec:136-q3-low-filtration} can have nonzero
graph residue even when every old degree-three Moore pair
has common-graph primitives.  The degree-two branch shows
that an immediate RH1 step does not rescue unrestricted
preservation either.  These statements do not refute a
construction retaining stronger data, choosing different old
maps, or maintaining vanishing old \(H^2\) before a
suitable special stage.  They do not construct a full common
second skeleton, and they do not decide integral connectedness.

\paragraph{Verification status.}
The author completed a whole saved-proof reading on 22 November
2026, including both exact low cohomology calculations, the
all-modulus primitive formulas, the actual two-truncation, the
free PD coordinate changes, and the Stokes obstruction.  Display
escapes and the explanation of the remaining nonunit odd block
were corrected during that reading.  The parent and reviewer1348
subsequently read the complete corrected saved proof and passed it.
They are recorded as contributor whole-proof readings: the parent
contributed the degree-two branch, and reviewer1348 contributed
the exact supplied-skeleton adaptation and audited the earlier
calculations.  Reviewer123's first complete saved-file reading
found the same corrected nonunit-block explanation and spacing
escapes; the targeted reread of those corrections completed a
whole-proof PASS.  That review also has contributor scope for
the mixed formula and degree-three coordinate discussion.
Earlier communicated calculations are not counted
as complete saved-file reviews.  The later minimal-order
refinement and retained-intermediate corollary have separate
whole-added-proof check events; a review of the original theorem
is not silently extended to those additions.  Model effort is
not exposed.
